\SourceLeaf{83}\emph{etiam ipsius $p^n$; qui vero ipsius $p$ est non-residuum, etiam ipsius $p^n$ non-residuum erit.}

Pars posterior huius propositionis per se est manifesta. Si itaque prior falsa esset, inter numeros ipso $p^n$ minores simulque per $p$ non divisibiles plures forent residua ipsius $p$ quam ipsius $p^n$, i. e. plures quam $\frac{1}{2}p^{n-1}(p-1)$. Nullo vero negotio perspici poterit, multitudinem residuorum numeri $p$ inter illos numeros esse praecise $=\frac{1}{2}p^{n-1}(p-1)$.

Aeque facile est, quadratum reipsa invenire, quod secundum modulum $p^n$ residuo dato sit congruum, si quadratum huic residuo secundum modulum $p$ congruum habetur.

Scilicet si quadratum habetur, $aa$, quod residuo dato $A$ secundum modulum $p^\mu$ est congruum, deducitur inde quadratum ipsi $A$ secundum modulum $p^\nu$ congruum (ubi $\nu>\mu$ et $=$ vel $<2\mu$ supponitur) sequenti modo. Ponatur radix quadrati quaesiti $=\pm a+xp^\mu$, quam formam eam habere debere facile perspicitur; debetque esse $aa\pm2axp^\mu+xxp^{2\mu}\equiv A\,(\mathrm{mod.}\,p^\nu)$ sive propter $2\mu\geq\nu$. $A-aa\equiv\pm2axp^\mu\,(\mathrm{mod.}\,p^\nu)$. Sit $A-aa=p^{\mu}d$, eritque $x$ valor expressionis $\pm\frac{d}{2a}\,(\mathrm{mod.}\,p^{\nu-\mu})$, quae huic $\pm\frac{A-aa}{2ap^\mu}\,(\mathrm{mod.}\,p^\nu)$ aequivalet.

Dato igitur quadrato ipsi $A$ secundum $p$ congruo, deducitur inde quadratum ipsi $A$ secundum modulum $p^2$ congruum; hinc ad modulum $p^4$, hinc ad $p^8$ etc. ascendi poterit.

\emph{Ex.} Proposito residuo $6$, quod secundum modulum $5$ quadrato $1$ congruum, invenitur quadratum $9^2$ cui secundum $25$ est congruum, $16^2$ cui secundum $125$. congruum etc.

\Article{102}

Quod vero attinet ad numeros per $p$ divisibiles, patet, eorum quadrata per $pp$ fore divisibilia, adeoque omnes numeros per $p$ quidem divisibiles, neque vero per $pp$, ipsius $p^n$ fore non residua. Generaliter vero, si proponitur numerus $p^k A$ ubi $A$ per $p$ non est divisibilis, hi casus erunt distinguendi:

1) Quando $k=$ vel $>n$, erit $p^k A\equiv0\,(\mathrm{mod.}\,p^n)$, i. e. residuum.

2) Quando $k<n$ atque impar. erit $p^k A$ non-residuum.

Si enim esset $p^k A=p^{2\kappa+1}A\equiv ss\,(\mathrm{mod.}\,p^n)$, $ss$ per $p^{2\kappa+1}$ divisibilis esset, id quod aliter fieri nequit, quam si fuerit $s$ per $p^{\kappa+1}$ divisibilis. Tunc vero $ss$
\SourceLeaf{84}etiam per $p^{2\kappa+2}$ divisibilis, adeoque etiam (quia $2\kappa+2$ certo non maior quam $n$) $p^k A$ \emph{i. e.} $p^{2\kappa+1}A$; sive $A$ per $p$, contra hyp.

3) Quando $k<n$ atque par. Tum $p^k A$ erit residuum vel non-residuum ipsius $p^n$, prout $A$ est residuum vel non-residuum ipsius $p$. Quando enim $A$ est residuum ipsius $p$, erit etiam residuum ipsius $p^{n-k}$. Posito autem $A\equiv aa (\mathrm{mod.}\,p^{n-k})$ erit $Ap^k\equiv aap^k\,(\mathrm{mod.}\,p^n)$, $aap^k$ vero est quadratum. Quando autem $A$ est non-residuum ipsius $p$, $p^k A$ residuum ipsius $p^n$ esse nequit. Ponatur enim $p^k A\equiv aa\,(\mathrm{mod.}\,p^n)$, eritque necessario $aa$ per $p^k$ divisibilis. Quotiens erit quadratum cui $A$ secundum modulum $p^{n-k}$ adeoque etiam secundum modulum $p$ congruus, \emph{i. e.} $A$ erit residuum ipsius $p$ contra hyp.

\Article{103}

Quoniam casum $p=2$ exclusimus, de hoc adhuc quaedam dicenda. Quando numerus $2$ est modulus, numerus quicunque erit residuum, non-residua nulla erunt. Quando vero $4$ est modulus, omnes numeri impares formae $4k+1$ erunt residua, omnes vero formae $4k+3$ non-residua. Tandem quando $8$ aut altior potestas numeri $2$ est modulus, omnes numeri impares formae $8k+1$ erunt residua, reliqui vero, seu ii qui sunt formarum $8k+3,\ 8k+5,\ 8k+7$, erunt non-residua. Pars posterior huius propositionis inde clara, quod quadratum cuiusvis numeri imparis, sive sit formae $4k+1$, sive formae $4k-1$, fit formae $8k+1$. Priorem ita probamus.

1) Si duorum numerorum vel summa vel differentia per $2^{n-1}$ est divisibilis, numerorum quadrata erunt congrua secundum modulum $2^n$. Si enim alter ponitur $=a$, erit alter formae $2^{n-1}h\pm a$, cuius quadratum invenitur $\equiv aa\,(\mathrm{mod.}\,2^n)$.

2) Quivis numerus impar, qui ipsius $2^n$ est residuum quadraticum, congruus erit quadrato alicui, cuius radix est numerus impar et $<2^{n-2}$. Sit enim quadratum quodcunque, cui numerus ille congruus, $aa$ atque numerus $a\equiv\pm\alpha (\mathrm{mod.}\,2^{n-1})$ ita ut $\alpha$ moduli semissem non superet (art. 4), eritque $aa\equiv\alpha\alpha$. Quare etiam numerus propositus erit $\equiv\alpha\alpha$. Manifesto vero tum $a$ tum $\alpha$ erunt impares atque $\alpha<2^{n-2}$.

3) Omnium numerorum imparium ipso $2^{n-2}$ minorum quadrata secundum $2^n$ incongrua erunt. Sint enim duo tales numeri $r$ et $s$, quorum quadrata si secundum $2^n$ essent congrua, foret $(r-s)(r+s)$ per $2^n$ divisibilis (posito $r>s$). Facile vero perspicitur numeros $r-s,\ r+s$ simul per $4$ divisibiles esse non
\SourceLeaf{85}posse, quare si alter tantummodo per $2$ est divisibilis, alter, ut productum per $2^n$ divisibilis fieret, per $2^{n-1}$ divisibilis esse deberet. \emph{Q. E. A.} quoniam uterque $<2^{n-2}$.

4) Quodsi denique haec quadrata ad \emph{residua} sua \emph{minima positiva} reducuntur, habebuntur $2^{n-3}$ residua quadratica diversa modulo minora\HistoricalNote{*)}{Puta quoniam multitudo numerorum imparium infra $2^{n-2}$ est $2^{n-3}$.}, quorum quodvis erit formae $8k+1$. Sed quum praecise $2^{n-3}$ numeri formae $8k+1$ modulo minores exstent, necessario hi omnes inter illa residua reperientur. \emph{Q. E. D.}

Ut quadratum numero dato formae $8k+1$ secundum modulum $2^n$ congruum inveniatur, methodus similis adhiberi potest, ut in art. 101; vid. etiam art. 88. --- Denique de numeris paribus eadem valent, quae art. 102 generaliter exposuimus.

\Article{104}

Circa multitudinem valorum diversorum (\emph{i. e.} secundum modulum incongruorum), quos expressio talis $V=\surd A\,(\mathrm{mod.}\,p^n)$ admittit, siquidem $A$ est residuum ipsius $p^n$, facile e praecc. colliguntur haec. (Numerum $p$ supponimus esse primum, ut ante, et brevitatis caussa casum $n=1$ statim includimus). I. Si $A$ per $p$ non est divisibilis, $V$ \emph{unum} valorem habet pro $p=2,\ n=1$, puta $V \equiv1$; \emph{duos}, quando $p$ est impar, nec non pro $p=2,\ n=2$, puta ponendo unum $\equiv v$, alter erit $\equiv-v$; \emph{quatuor} pro $p=2,\ n>2$. scilicet ponendo unum $\equiv v$, reliqui erunt $\equiv-v,\ 2^{n-1}+v,\ 2^{n-1}-v$. II. Si $A$ per $p$ divisibilis est, neque vero per $p^n$, sit potestas altissima ipsius $p$ ipsum $A$ metiens $p^{2\mu}$ (manifesto enim ipsius exponens par esse debebit) atque $A=ap^{2\mu}$. Tunc patet, omnes valores ipsius $V$ per $p^\mu$ divisibiles esse, et quotientes e divisione ortos fieri valores expr. $V'=\surd a\,(\mathrm{mod.}\,p^{n-2\mu})$; hinc omnes valores diversi ipsius $V$ prodibunt, multiplicando omnes valores expr. $V'$ inter $0$ et $p^{n-\mu}$ sitos per $p^\mu$; quare illi exhibebuntur per

\[
vp^\mu,\quad
vp^\mu+p^{n-\mu},\quad
vp^\mu+2p^{n-\mu}\ldots vp^\mu+(p^\mu-1)p^{n-\mu}
\]

si $v$ indefinite omnes valores \emph{diversos} expr. $V'$ exprimit, ita ut illorum multitudo fiat $p^\mu,\ 2p^\mu$ vel $4p^\mu$, prout multitudo horum (per casum I) est $1,\ 2$ vel $4$. III. Si $A$ per $p^n$ divisibilis est, facile perspicitur, statuendo $n=2m$ vel $=2m-1$, prout par est vel impar, omnes numeros per $p^m$ divisibiles, neque ullos alios, esse valores ipsius $V$; quare omnes valores diversi hi erunt $0,p^m,2p^m\ldots(p^{n-m}-1)p^m$, quorum multitudo $p^{n-m}$.

\SourceLeaf{86}\Article{105}

Superest casus. ubi modulus $m$ e pluribus numeris primis compositus est. Sit $m=abc\ldots$, designantibus $a,b,c$ etc. numeros primos diversos aut primorum diversorum potestates, patetque statim, si $n$ sit residuum ipsius $m$, fore etiam $n$ residuum singulorum $a,b,c$ etc., adeoque $n$ certo non-residuum ipsius $m$ esse, si fuerit NR. ullius e numeris $a,b,c$ etc. Vice versa autem, si $n$ singulorum $a,b,c$ etc. residuum est, etiam residuum producti $m$ erit. Supponendo enim, $n\equiv A^2$, $B^2,C^2$ etc. sec. mod. $a,b,c$ etc. resp., patet, si numerus $N$ ipsis $A,B,C$ etc. sec. mod. $a,b,c$ etc. resp. congruus eruatur (art. 32), fore $n\equiv NN$ secundum omnes hos modulos adeoque etiam secundum productum $m$. --- Quum facile perspiciatur, hoc modo e combinatione \emph{cuiusvis} valoris ipsius $A$ sive expr. $\surd n\,(\mathrm{mod.}\,a)$ cum \emph{quovis} valore ipsius $B$ cum \emph{quovis} valore ipsius $C$ etc. oriri valorem ipsius $N$ sive expr. $\surd n\,(\mathrm{mod.}\,m)$, nec non e combinationibus diversis produci diversos $N$, et e cunctis cunctos: multitudo omnium valorum diversorum ipsius $N$ aequalis erit producto e multitudinibus valorum ipsorum $A,B,C$ etc. quas determinare in art. praec. docuimus. --- Porro manifestum est, si unus valor expressionis $\surd n\,(\mathrm{mod.}\,m)$ sive ipsius $N$ fuerit notus, hunc simul fore valorem omnium $A,B,C$ etc.; et quum hinc per art. praec. omnes reliqui valores harum quantitatum deduci possint, facile sequitur, ex uno valore ipsius $N$ omnes reliquos obtineri posse.

\emph{Ex.} Sit modulus $315$ cuius residuum an non-residuum sit $46$, quaeritur. Divisores primi numeri $315$ sunt $3,5,7$, atque numerus $46$ residuum cuiusvis eorum quare etiam ipsius $315$ erit residuum. Porro. quia $46\equiv1$, et $\equiv64 (\mathrm{mod.}\,9)$; $\equiv1$ et $\equiv16\,(\mathrm{mod.}\,5)$; $\equiv4$ et $\equiv25\,(\mathrm{mod.}\,7)$, inveniuntur radices quadratorum, quibus $46$ secundum modulum $315$ congruus. $19,26,44,89,226$, $271,289,296$.

\begin{center}
\small\itshape
Criterium generale, utrum numerus datus numeri primi dati residuum sit an non-residuum.
\end{center}

\Article{106}

Ex praecedentibus colligitur, si tantummodo semper dignosci possit utrum \emph{numerus primus} datus numeri \emph{primi dati} residuum sit an non-residuum, omnes reliquos casus ad hunc reduci posse. Pro illo itaque casu criteria certa omni studio nobis erunt indaganda. Antequam autem hanc perquisitionem aggrediamur, criterium quoddam exhibemus ex Sect. praec. petitum, quod quamvis in praxi nul%
